
\documentclass{article}
\usepackage{colortbl}
\usepackage{makecell}
\usepackage{multirow}
\usepackage{supertabular}

\begin{document}

\newcounter{utterance}

\centering \large Interaction Transcript for game `cladder', experiment `full\_v1.5\_default', episode 4010 with llama{-}3.1{-}8b{-}instruct.
\vspace{24pt}

{ \footnotesize  \setcounter{utterance}{1}
\setlength{\tabcolsep}{0pt}
\begin{supertabular}{c@{$\;$}|p{.15\linewidth}@{}p{.15\linewidth}p{.15\linewidth}p{.15\linewidth}p{.15\linewidth}p{.15\linewidth}}
    \# & \multicolumn{2}{c}{Player} && \multicolumn{2}{c}{Game Master} \\
    \hline

    \theutterance \stepcounter{utterance}  
    & & & \multicolumn{4}{p{0.6\linewidth}}{
        \cellcolor[rgb]{0.9,0.9,0.9}{
            \makecell[{{p{\linewidth}}}]{
                \texttt{\tiny{[P1$\langle$GM]}}
                \texttt{You are an expert in causal inference. The following question is not a typical commonsense query, but rather a meticulously designed question created by a professor specializing in causal inference, intended to assess the students' mastery of the course content.} \\
\texttt{Imagine a self{-}contained, hypothetical world with only the following conditions, and without any unmentioned factors or causal relationships: Liking spicy food has a direct effect on parents' social status and child's intelligence. Other unobserved factors has a direct effect on parents' social status and child's intelligence. Parents' social status has a direct effect on child's intelligence. Other unobserved factors is unobserved. The overall probability of liking spicy food is 51\%. The probability of not liking spicy food and intelligent child is 18\%. The probability of liking spicy food and intelligent child is 34\%. Is the chance of intelligent child larger when observing liking spicy food?} \\
\texttt{Start your answer with 'Yes' or 'No', followed by additional reasoning or evidence to support your explanation.} \\
            }
        }
    }
     \\ \\

    \theutterance \stepcounter{utterance}  
    & \multicolumn{4}{p{0.6\linewidth}}{
        \cellcolor[rgb]{0.561,0.737,0.561}{
            \makecell[{{p{\linewidth}}}]{
                \texttt{\tiny{[P1$\rangle$GM]}}
                \texttt{No.} \\
\\ 
\texttt{To answer this question, we need to apply the principles of causal inference, specifically the concept of confounding variables and the problem of selection bias. } \\
\\ 
\texttt{Let's denote the variables as follows:} \\
\texttt{{-} L: Liking spicy food} \\
\texttt{{-} S: Parents' social status} \\
\texttt{{-} I: Child's intelligence} \\
\texttt{{-} U: Other unobserved factors (which affects S and I)} \\
\\ 
\texttt{We are given the following information:} \\
\texttt{{-} P(L) = 0.51 (the probability of liking spicy food)} \\
\texttt{{-} P(¬L, I) = 0.18 (the probability of not liking spicy food and intelligent child)} \\
\texttt{{-} P(L, I) = 0.34 (the probability of liking spicy food and intelligent child)} \\
\\ 
\texttt{We want to determine if the chance of an intelligent child is larger when observing liking spicy food, i.e., whether P(I|L) > P(I|¬L).} \\
\\ 
\texttt{To do this, we can use the law of total probability and the concept of conditional probability. We can express P(I) as follows:} \\
\\ 
\texttt{P(I) = P(I|L)P(L) + P(I|¬L)P(¬L)} \\
\\ 
\texttt{We can also express P(I|L) and P(I|¬L) using the law of total probability:} \\
\\ 
\texttt{P(I|L) = P(I, L) / P(L)} \\
\texttt{P(I|¬L) = P(I, ¬L) / P(¬L)} \\
\\ 
\texttt{Substituting the given values, we get:} \\
\\ 
\texttt{P(I) = (0.34 / 0.51)P(L) + (0.18 / P(¬L))} \\
\\ 
\texttt{Since P(¬L) = 1 {-} P(L) = 0.49, we can substitute this value:} \\
\\ 
\texttt{P(I) = (0.34 / 0.51)P(L) + (0.18 / 0.49)} \\
\\ 
\texttt{Now, we want to compare P(I|L) and P(I|¬L). We can do this by comparing the coefficients of P(L) and P(¬L) in the expression for P(I):} \\
\\ 
\texttt{P(I|L) = 0.34 / 0.51} \\
\texttt{P(I|¬L) = 0.18 / 0.49} \\
\\ 
\texttt{To determine if P(I|L) > P(I|¬L), we can compare the ratios:} \\
\\ 
\texttt{(0.34 / 0.51) / (0.18 / 0.49) = (0.34 * 0.49) / (0.18 * 0.51) = 0.1666 / 0.0922 ≈ 1.81} \\
\\ 
\texttt{Since the ratio is greater than 1, we can conclude that P(I|L) > P(I|¬L). However, this is not the case. The ratio is actually less than 1, which means that P(I|L) < P(I|¬L).} \\
\\ 
\texttt{However} \\
            }
        }
    }
    & & \\ \\

    \theutterance \stepcounter{utterance}  
    & & & \multicolumn{2}{p{0.3\linewidth}}{
        \cellcolor[rgb]{0.9,0.9,0.9}{
            \makecell[{{p{\linewidth}}}]{
                \texttt{\tiny{[GM$|$GM]}}
                \texttt{No} \\
            }
        }
    }
    & & \\ \\

    \theutterance \stepcounter{utterance}  
    & & & \multicolumn{2}{p{0.3\linewidth}}{
        \cellcolor[rgb]{0.9,0.9,0.9}{
            \makecell[{{p{\linewidth}}}]{
                \texttt{\tiny{[GM$|$GM]}}
                \texttt{yes} \\
            }
        }
    }
    & & \\ \\

    \theutterance \stepcounter{utterance}  
    & & & \multicolumn{2}{p{0.3\linewidth}}{
        \cellcolor[rgb]{0.9,0.9,0.9}{
            \makecell[{{p{\linewidth}}}]{
                \texttt{\tiny{[GM$|$GM]}}
                \texttt{game\_result = LOSE} \\
            }
        }
    }
    & & \\ \\

\end{supertabular}
}

\end{document}
